\documentclass[aspectratio=169, 10pt]{beamer}

% --- Packages fondamentaux ---
\usepackage[utf8]{inputenc}
\usepackage[T1]{fontenc}
\usepackage[french]{babel}
\usepackage{lmodern}
\usepackage{microtype}
\usepackage{booktabs}
\usepackage{colortbl}
\usepackage{amssymb}
\usepackage{tikz}
\usepackage{xcolor}
\usepackage{hyperref}

% --- Bibliothèques TikZ ---
\usetikzlibrary{shapes.geometric, arrows.meta, positioning, calc, shadows, fit}

% --- Palette Corporate Capgemini & Assurance Qualité ---
\definecolor{CapgeminiBlue}{RGB}{0, 112, 173}       % Bleu officiel Capgemini #0070AD
\definecolor{CapgeminiDark}{RGB}{11, 25, 44}        % Bleu nuit #0B192C
\definecolor{CapgeminiVibrant}{RGB}{0, 163, 224}    % Cyan dynamique #00A3E0
\definecolor{TestSuccess}{RGB}{16, 185, 129}        % Vert validation #10B981
\definecolor{TestWarning}{RGB}{245, 158, 11}        % Orange attention #F59E0B
\definecolor{CardBg}{RGB}{248, 250, 252}            % Fond carte #F8FAFC
\definecolor{BorderColor}{RGB}{226, 232, 240}       % Bordure #E2E8F0
\definecolor{TextPrimary}{RGB}{15, 23, 42}          % Texte foncé #0F172A
\definecolor{TextSecondary}{RGB}{100, 116, 139}     % Texte gris #64748B

% --- Configuration des marges Beamer standards (pour éviter les débordements) ---
\setbeamersize{text margin left=0.8cm, text margin right=0.8cm}
\setbeamertemplate{navigation symbols}{}
\setbeamercolor{background canvas}{bg=white}
\setbeamercolor{normal text}{fg=TextPrimary}

% --- En-tête propre sans débordement ---
\setbeamertemplate{frametitle}{
    \vspace*{0.35cm}
    \begin{minipage}{\textwidth}
        {\color{CapgeminiBlue}\bfseries\large\insertframetitle}
        \ifx\insertframesubtitle\@empty\else
            \\[0.06cm]{\color{TextSecondary}\scriptsize\insertframesubtitle}
        \fi
        \\[0.12cm]
        {\color{CapgeminiBlue}\rule{2.8cm}{1.6pt}}
    \end{minipage}
    \vspace*{-0.15cm}
}

% --- Pied de page ---
\setbeamertemplate{footline}{
    \begin{beamercolorbox}[wd=\paperwidth, ht=0.55cm, dp=0.2cm, leftskip=0.8cm, rightskip=0.8cm]{footline}
        \begin{tikzpicture}[remember picture, overlay]
            \draw[BorderColor, line width=0.5pt] ($(current page.south west)+(0.8cm,0.55cm)$) -- ($(current page.south east)+(-0.8cm,0.55cm)$);
            \node[anchor=south west, xshift=0.8cm, yshift=0.15cm] at (current page.south west) {
                \tiny\color{TextSecondary}\textbf{Rafik AIT ICHOU} \ \textbar\ \ Stage PFE QA \& Test Logiciel \ \textbar\ \ \textbf{Capgemini}
            };
            \node[anchor=south east, xshift=-0.8cm, yshift=0.15cm] at (current page.south east) {
                \tiny\color{CapgeminiBlue}\textbf{\insertframenumber} \color{TextSecondary}/ \inserttotalframenumber
            };
        \end{tikzpicture}
    \end{beamercolorbox}
}

% --- Style de boîte carte ---
\tikzset{
    card/.style={
        rectangle, rounded corners=5pt, draw=BorderColor, fill=CardBg,
        inner sep=6pt, text=TextPrimary
    },
    chip/.style={
        rectangle, rounded corners=3pt, fill=CapgeminiBlue!10, draw=CapgeminiBlue!40,
        inner sep=3pt, font=\scriptsize\bfseries, text=CapgeminiBlue
    }
}

% ==============================================================================
% DÉBUT DU DOCUMENT
% ==============================================================================
\begin{document}

% ------------------------------------------------------------------------------
% DIAPOSITIVE 1 : Titre & Introduction
% ------------------------------------------------------------------------------
\begin{frame}[plain]
    \begin{tikzpicture}[remember picture, overlay]
        % Fond sombre exécutif
        \fill[CapgeminiDark] (current page.north west) rectangle (current page.south east);
        
        % Cercles décoratifs Capgemini
        \fill[CapgeminiBlue, opacity=0.85] ($(current page.north east)+(-2.8cm, 0.5cm)$) circle (4cm);
        \fill[CapgeminiVibrant, opacity=0.35] ($(current page.south east)+(-1.2cm, -0.5cm)$) circle (3cm);

        % Contenu centré verticalement
        \node[anchor=west, xshift=1cm, yshift=0.1cm] at (current page.west) {
            \begin{minipage}{0.78\paperwidth}
                % Badge haut
                \begin{tikzpicture}
                    \node[rectangle, rounded corners=3pt, fill=CapgeminiVibrant!20, draw=CapgeminiVibrant, text=white, inner sep=4pt] {
                        \scriptsize\bfseries CANDIDATURE STAGE DE FIN D'ÉTUDES (PFE 2027) \textbullet\ CAPGEMINI
                    };
                \end{tikzpicture}
                \vspace{0.35cm}
                
                % Titre principal
                {\fontsize{20}{24}\selectfont \color{white}\bfseries Ingénierie QA, Automatisation\\[0.1cm] des Tests \& Pratiques Scrum}\\[0.3cm]
                {\fontsize{10.5}{13}\selectfont \color{CapgeminiVibrant}\textbf{Assurer la qualité, la fiabilité et la sécurité des logiciels critiques}}\\[0.45cm]
                
                % Barre de séparation
                {\color{CapgeminiVibrant}\rule{4cm}{1.8pt}}\\[0.35cm]
                
                % Profil candidat
                {\large\color{white}\textbf{Rafik AIT ICHOU}}\\[0.08cm]
                {\small\color{gray!30} Élève-Ingénieur en Informatique \& Réseaux — \textbf{EMSI Marrakech}}\\[0.06cm]
                {\scriptsize\color{gray!40} Disponible pour stage de 6 mois dès le \textbf{25 Janvier 2027} \textbullet\ Mobilité : Casablanca, Rabat, International}
                \vspace{0.45cm}

                % Boîte métriques clés (Fond foncé, texte clair et lisible)
                \begin{tikzpicture}
                    \node[rectangle, rounded corners=4pt, fill=white!12, draw=CapgeminiVibrant!50, inner sep=5pt] {
                        \footnotesize\color{white}
                        \textbf{380+} Tests Automatisés \ \textbar\ \ 
                        \textbf{86.8\%} Couverture SonarQube (OCP Safi) \ \textbar\ \ 
                        \textbf{Scrum \& ISTQB} Mindset
                    };
                \end{tikzpicture}
            \end{minipage}
        };
    \end{tikzpicture}
\end{frame}

% ------------------------------------------------------------------------------
% DIAPOSITIVE 2 : Philosophie QA & Pyramide des Tests
% ------------------------------------------------------------------------------
\begin{frame}{Philosophie QA \& Stratégie de Test}{Approche Shift-Left : tester tôt pour corriger vite et à moindre coût}
    \begin{columns}[c]
        % Colonne Gauche : Schéma TikZ Pyramide bien proportionnée
        \begin{column}{0.48\textwidth}
            \centering
            \begin{tikzpicture}[scale=0.82]
                % Étage 3 (Sommet) : Tests E2E & Sécurité
                \draw[fill=CapgeminiDark, draw=CapgeminiDark] (-1.0, 2.7) -- (1.0, 2.7) -- (0, 3.8) -- cycle;
                \node at (0, 3.15) {\color{white}\bfseries\scriptsize Tests E2E \& Sécurité};
                \node at (0, 2.85) {\color{CapgeminiVibrant}\fontsize{5.5}{6.5}\selectfont Selenium \textbullet\ OWASP \textbullet\ JMeter};

                % Étage 2 (Milieu) : Tests d'Intégration & API
                \draw[fill=CapgeminiBlue, draw=CapgeminiBlue] (-1.9, 1.4) -- (1.9, 1.4) -- (1.05, 2.65) -- (-1.05, 2.65) -- cycle;
                \node at (0, 2.1) {\color{white}\bfseries\scriptsize Tests d'Intégration \& API};
                \node at (0, 1.65) {\color{white!90}\fontsize{5.8}{7}\selectfont Postman \textbullet\ Spring Boot Test \textbullet\ SQL};

                % Étage 1 (Base) : Tests Unitaires
                \draw[fill=TestSuccess, draw=TestSuccess] (-2.7, 0.0) -- (2.7, 0.0) -- (1.95, 1.35) -- (-1.95, 1.35) -- cycle;
                \node at (0, 0.8) {\color{white}\bfseries\small Tests Unitaires \& Mocking};
                \node at (0, 0.35) {\color{white!95}\fontsize{6}{7}\selectfont JUnit 5 \textbullet\ Mockito \textbullet\ Pest (86.8\% couverture)};

                % Flèche gauche : Coût
                \draw[-{Stealth[length=2mm]}, thick, CapgeminiBlue] (-3.1, 0.1) -- (-3.1, 3.6)
                    node[midway, left=1pt, rotate=90, font=\tiny\bfseries, color=CapgeminiBlue] {Coût et Temps de test};
                
                % Flèche droite : Vélocité
                \draw[-{Stealth[length=2mm]}, thick, TestSuccess] (3.1, 3.6) -- (3.1, 0.1)
                    node[midway, right=1pt, rotate=-90, font=\tiny\bfseries, color=TestSuccess] {Rapidité et Fréquence};
            \end{tikzpicture}
        \end{column}

        % Colonne Droite : Explications simples et claires
        \begin{column}{0.52\textwidth}
            \begin{tikzpicture}
                \node[card, text width=0.94\linewidth] {
                    \textbf{\color{CapgeminiBlue}\footnotesize 1. Le principe du Shift-Left}
                    \vspace{0.08cm}
                    \begin{itemize}\setlength{\itemsep}{2pt}\scriptsize
                        \item Définir les tests \textbf{dès la rédaction des User Stories}.
                        \item Automatiser pendant le sprint, pas à la fin du projet.
                        \item Corriger les défauts immédiatement quand c'est simple.
                    \end{itemize}
                };
            \end{tikzpicture}
            \vspace{0.15cm}

            \begin{tikzpicture}
                \node[card, text width=0.94\linewidth] {
                    \textbf{\color{CapgeminiBlue}\footnotesize 2. Les règles de qualité appliquées}
                    \vspace{0.08cm}
                    \begin{itemize}\setlength{\itemsep}{2pt}\scriptsize
                        \item \textbf{Priorité par le risque} : tester d'abord les flux critiques.
                        \item \textbf{Non-régression} : tests rejoués à chaque modification.
                        \item \textbf{Quality Gate} : aucun code déployé sans validation SonarQube.
                    \end{itemize}
                };
            \end{tikzpicture}
        \end{column}
    \end{columns}
\end{frame}

% ------------------------------------------------------------------------------
% DIAPOSITIVE 3 : Projet Phare — Smart HSE 360 (OCP Safi)
% ------------------------------------------------------------------------------
\begin{frame}{Projet Phare : Smart HSE 360 (OCP Safi)}{Plateforme industrielle de gestion des incidents HSE \textbullet\ Automatisation \& Quality Gate}
    \begin{columns}[c]
        % Colonne Gauche : Rôle et réalisations concrètes
        \begin{column}{0.50\textwidth}
            \begin{tikzpicture}
                \node[card, text width=0.94\linewidth] {
                    \textbf{\color{CapgeminiBlue}\small Mon Rôle \& Actions QA :}
                    \vspace{0.1cm}
                    \begin{itemize}\setlength{\itemsep}{3pt}\scriptsize
                        \item \textbf{Plan de test complet} : Validation de l'API Spring Boot, du web Angular et de l'app mobile Flutter.
                        \item \textbf{164 Tests Unitaires (JUnit 5, Mockito)} : Vérification des calculs de gravité et des règles de notification.
                        \item \textbf{58 Tests E2E (Selenium WebDriver)} : Parcours complet de la déclaration d'incident à sa clôture.
                        \item \textbf{Quality Gate SonarQube} : 0 bug bloquant, 0 faille de sécurité, couverture validée à \textbf{86.8\%}.
                    \end{itemize}
                };
            \end{tikzpicture}
            \vspace{0.15cm}

            \begin{tikzpicture}
                \node[card, fill=TestSuccess!8, draw=TestSuccess!40, text width=0.94\linewidth] {
                    \centering
                    \scriptsize
                    \textbf{\color{TestSuccess}Résultat mesurable :} 222 tests passants en intégration continue\\
                    \textbf{Attestation officielle de satisfaction délivrée par l'OCP Safi}.
                };
            \end{tikzpicture}
        \end{column}

        % Colonne Droite : Pipeline CI/CD simple et clair
        \begin{column}{0.50\textwidth}
            \centering
            {\footnotesize\textbf{\color{CapgeminiDark}Pipeline de Validation Déployé}}\\[0.15cm]
            \begin{tikzpicture}[
                node distance=0.28cm,
                pipe/.style={rectangle, rounded corners=3pt, draw=BorderColor, fill=white, minimum width=4.8cm, minimum height=0.48cm, font=\scriptsize, align=center}
            ]
                \node[pipe, fill=gray!8] (p1) {\textbf{1. Git Push / PR} \color{TextSecondary}(Branche de travail)};
                \node[pipe, below=of p1, fill=CapgeminiBlue!8, draw=CapgeminiBlue!30] (p2) {\textbf{2. Build Maven} \color{CapgeminiBlue}(Compilation et Lint)};
                \node[pipe, below=of p2, fill=TestSuccess!12, draw=TestSuccess!40] (p3) {\textbf{3. Tests Unitaires} \color{TextPrimary}(164 tests JUnit 5 / Mockito)};
                \node[pipe, below=of p3, fill=CapgeminiVibrant!12, draw=CapgeminiVibrant!40] (p4) {\textbf{4. Tests E2E Selenium} \color{TextPrimary}(58 scénarios d'usage)};
                \node[pipe, below=of p4, fill=TestWarning!15, draw=TestWarning!50] (p5) {\textbf{5. SonarQube Gate} \color{TextPrimary}(Couverture $\ge 85\%$)};
                \node[pipe, below=of p5, fill=CapgeminiDark, draw=CapgeminiDark] (p6) {\color{white}\textbf{6. Déploiement Docker Staging}};

                \draw[-{Stealth[length=1.8mm]}, thick, CapgeminiBlue] (p1) -- (p2);
                \draw[-{Stealth[length=1.8mm]}, thick, CapgeminiBlue] (p2) -- (p3);
                \draw[-{Stealth[length=1.8mm]}, thick, TestSuccess] (p3) -- (p4);
                \draw[-{Stealth[length=1.8mm]}, thick, CapgeminiVibrant] (p4) -- (p5);
                \draw[-{Stealth[length=1.8mm]}, thick, TestWarning] (p5) -- (p6);
            \end{tikzpicture}
        \end{column}
    \end{columns}
\end{frame}

% ------------------------------------------------------------------------------
% DIAPOSITIVE 4 : Systèmes Connectés, IoT & Stage ENS Marrakech
% ------------------------------------------------------------------------------
\begin{frame}{Tests IoT, Systèmes Connectés \& Stage ENS}{Validation matériel-logiciel : vision par ordinateur, capteurs et tests d'intégration}
    \begin{columns}[c]
        % Carte Projet 1 : Smart Parking System
        \begin{column}{0.49\textwidth}
            \begin{tikzpicture}
                \node[card, text width=0.94\linewidth] {
                    \textbf{\color{CapgeminiBlue}\small Smart Parking (IoT, Vision \& Next.js)}
                    \vspace{0.08cm}
                    \begin{itemize}\setlength{\itemsep}{2pt}\scriptsize
                        \item \textbf{Électronique \& Capteurs} : Caméras ESP32-CAM et Raspberry Pi exécutant des modèles OpenCV / YOLOv8.
                        \item \textbf{Tests d'intégration IoT} : validation des flux WebSockets et MQTT en temps réel (latence $< 500$ ms).
                        \item \textbf{Tests unitaires \& E2E} : logique de réservation sous JUnit 5 et parcours de paiement simulés avec Selenium.
                        \item \textbf{Résilience réseau} : tests de déconnexion temporaire des caméras sans perte d'état.
                    \end{itemize}
                    \vspace{0.05cm}
                    \centering
                    \tikz{\node[chip, fill=CapgeminiBlue!10, text=CapgeminiBlue, draw=CapgeminiBlue!30]{IoT ESP32 \textbullet\ YOLOv8 \textbullet\ Tests Selenium};}
                };
            \end{tikzpicture}
        \end{column}

        % Carte Projet 2 : absENS (Stage 6 mois ENS Marrakech)
        \begin{column}{0.49\textwidth}
            \begin{tikzpicture}
                \node[card, text width=0.94\linewidth] {
                    \textbf{\color{CapgeminiBlue}\small absENS — Stage de 6 Mois (ENS Marrakech)}
                    \vspace{0.08cm}
                    \begin{itemize}\setlength{\itemsep}{2pt}\scriptsize
                        \item \textbf{Borne IoT \& Vision} : Module Raspberry Pi avec caméra OpenCV scannant des QR codes rotatifs chiffrés.
                        \item \textbf{Tests anti-fraude} : validation de l'impossibilité d'émarger par photo d'écran (durée de validité QR $< 3$ s).
                        \item \textbf{Tests de charge amphi} : émargement de 1\,200+ étudiants sans saturation de l'API Laravel / MySQL.
                        \item \textbf{Livrable académique} : attestation officielle de stage de 6 mois délivrée par l'ENS Marrakech.
                    \end{itemize}
                    \vspace{0.05cm}
                    \centering
                    \tikz{\node[chip, fill=TestSuccess!10, text=TestSuccess, draw=TestSuccess!30]{Stage 6 mois validé \textbullet\ Scan $< 3$s \textbullet\ 0\% Fraude};}
                };
            \end{tikzpicture}
        \end{column}
    \end{columns}
\end{frame}

% ------------------------------------------------------------------------------
% DIAPOSITIVE 5 : La QA au Cœur de l'Agilité Scrum
% ------------------------------------------------------------------------------
\begin{frame}{La QA au Cœur de l'Agilité Scrum}{Le rôle actif du testeur dans chaque étape du sprint}
    \centering
    \begin{tikzpicture}[
        scale=0.88,
        stepbox/.style={rectangle, rounded corners=5pt, draw=BorderColor, fill=white, minimum width=3.3cm, minimum height=3.3cm, inner sep=5pt}
    ]
        % Étape 1 : Planning
        \node[stepbox] (s1) at (0,0) {
            \begin{minipage}{3.1cm}
                \centering
                {\color{CapgeminiBlue}\textbf{\footnotesize 1. Sprint Planning}}\\[0.1cm]
                \begin{itemize}\setlength{\itemsep}{1pt}\tiny\raggedright
                    \item Écrire les critères d'acceptation (\textit{Given-When-Then}).
                    \item Estimer l'effort de test avec les développeurs.
                    \item Vérifier la \textbf{Definition of Ready (DoR)}.
                \end{itemize}
            \end{minipage}
        };

        % Étape 2 : Développement & Test
        \node[stepbox, right=0.35cm of s1] (s2) {
            \begin{minipage}{3.1cm}
                \centering
                {\color{CapgeminiBlue}\textbf{\footnotesize 2. En Cours de Sprint}}\\[0.1cm]
                \begin{itemize}\setlength{\itemsep}{1pt}\tiny\raggedright
                    \item Écrire les tests unitaires et d'API en parallèle du dev.
                    \item Tester les endpoints dès qu'ils sont prêts (Postman).
                    \item Partager les points de blocage au \textbf{Daily Scrum}.
                \end{itemize}
            \end{minipage}
        };

        % Étape 3 : Review & DoD
        \node[stepbox, right=0.35cm of s2] (s3) {
            \begin{minipage}{3.1cm}
                \centering
                {\color{TestSuccess}\textbf{\footnotesize 3. Sprint Review \& DoD}}\\[0.1cm]
                \begin{itemize}\setlength{\itemsep}{1pt}\tiny\raggedright
                    \item Vérifier la \textbf{Definition of Done (DoD)} :
                    \begin{itemize}\tiny
                        \item Tests unitaires passants à 100\%.
                        \item Couverture $\ge 80\%$.
                        \item Zéro anomalie bloquante.
                    \end{itemize}
                    \item Valider la démonstration client.
                \end{itemize}
            \end{minipage}
        };

        % Étape 4 : Rétrospective
        \node[stepbox, right=0.35cm of s3] (s4) {
            \begin{minipage}{3.1cm}
                \centering
                {\color{TestWarning}\textbf{\footnotesize 4. Rétrospective}}\\[0.1cm]
                \begin{itemize}\setlength{\itemsep}{1pt}\tiny\raggedright
                    \item Analyser les bugs trouvés pour comprendre leur cause.
                    \item Améliorer la suite de non-régression.
                    \item Optimiser le temps d'exécution des tests.
                \end{itemize}
            \end{minipage}
        };

        % Flèches de flux
        \draw[-{Stealth[length=2mm]}, line width=1.3pt, CapgeminiBlue] (s1.east) -- (s2.west);
        \draw[-{Stealth[length=2mm]}, line width=1.3pt, CapgeminiBlue] (s2.east) -- (s3.west);
        \draw[-{Stealth[length=2mm]}, line width=1.3pt, TestSuccess] (s3.east) -- (s4.west);
        \draw[-{Stealth[length=2mm]}, dashed, line width=1pt, TextSecondary] 
            ($(s4.south)+(0,0.2)$) -- ++(0,-0.3) -| ($(s1.south)+(0,0.2)$)
            node[pos=0.25, below, font=\tiny\bfseries, color=TextSecondary] {Amélioration Continue};
    \end{tikzpicture}
    
    \vspace{0.2cm}
    \begin{tikzpicture}
        \node[card, fill=CapgeminiBlue!5, draw=CapgeminiBlue!20, text width=0.96\linewidth] {
            \scriptsize
            \textbf{\color{CapgeminiBlue}Expérience Scrum (Président Club E-TEAM ENS) :} Animation des rituels agiles, coordination d'équipes de projets et culture du résultat testé et validé.
        };
    \end{tikzpicture}
\end{frame}

% ------------------------------------------------------------------------------
% DIAPOSITIVE 6 : Boîte à Outils QA & Stack Technologique
% ------------------------------------------------------------------------------
\begin{frame}{Boîte à Outils QA \& Stack Technologique}{Maîtrise pratique des outils de test et d'automatisation}
    \begin{columns}[c]
        % Quadrant 1 : Tests Fonctionnels & Qualité
        \begin{column}{0.49\textwidth}
            \begin{tikzpicture}
                \node[card, text width=0.93\linewidth] {
                    \textbf{\color{CapgeminiBlue}\footnotesize $\blacktriangleright$ Tests Fonctionnels \& API}
                    \vspace{0.06cm}
                    \begin{itemize}\setlength{\itemsep}{1.5pt}\scriptsize
                        \item \textbf{Selenium WebDriver} : Automatisation web (Page Object Model).
                        \item \textbf{Postman \& Newman} : Tests d'API REST et validation des réponses HTTP.
                        \item \textbf{Apache JMeter} : Tests de charge et mesure des temps de réponse.
                    \end{itemize}
                };
            \end{tikzpicture}
            \vspace{0.12cm}

            \begin{tikzpicture}
                \node[card, text width=0.93\linewidth] {
                    \textbf{\color{CapgeminiBlue}\footnotesize $\blacktriangleright$ Qualité de Code \& Analyse}
                    \vspace{0.06cm}
                    \begin{itemize}\setlength{\itemsep}{1.5pt}\scriptsize
                        \item \textbf{SonarQube} : Détection des bugs, failles de sécurité et code smells.
                        \item \textbf{JaCoCo} : Mesure exacte du taux de couverture de code.
                        \item \textbf{Clean Code} : Refactoring sûr et sans régression.
                    \end{itemize}
                };
            \end{tikzpicture}
        \end{column}

        % Quadrant 2 : Tests Unitaires & CI/CD
        \begin{column}{0.49\textwidth}
            \begin{tikzpicture}
                \node[card, text width=0.93\linewidth] {
                    \textbf{\color{CapgeminiBlue}\footnotesize $\blacktriangleright$ Tests Unitaires \& Mocking}
                    \vspace{0.06cm}
                    \begin{itemize}\setlength{\itemsep}{1.5pt}\scriptsize
                        \item \textbf{JUnit 5 / AssertJ} : Tests unitaires, assertions claires et complètes.
                        \item \textbf{Mockito} : Isolation des composants par simulation des dépendances.
                        \item \textbf{Pest / PHPUnit} : Tests rapides de logique métier.
                    \end{itemize}
                };
            \end{tikzpicture}
            \vspace{0.12cm}

            \begin{tikzpicture}
                \node[card, text width=0.93\linewidth] {
                    \textbf{\color{CapgeminiBlue}\footnotesize $\blacktriangleright$ Intégration Continue \& Environnement}
                    \vspace{0.06cm}
                    \begin{itemize}\setlength{\itemsep}{1.5pt}\scriptsize
                        \item \textbf{GitHub Actions} : Exécution automatique des tests à chaque PR.
                        \item \textbf{Docker} : Environnements de test identiques et reproductibles.
                        \item \textbf{Maven \& Git} : Gestion propre des dépendances et versions.
                    \end{itemize}
                };
            \end{tikzpicture}
        \end{column}
    \end{columns}
\end{frame}

% ------------------------------------------------------------------------------
% DIAPOSITIVE 7 : Pourquoi Rafik chez Capgemini ?
% ------------------------------------------------------------------------------
\begin{frame}{Pourquoi Capgemini \& Quelle Valeur Immédiate ?}{Motivation, rigueur technique et disponibilité pour le stage PFE}
    \begin{columns}[c]
        % Colonne 1 : Atouts pour Capgemini
        \begin{column}{0.58\textwidth}
            \begin{tikzpicture}
                \node[card, text width=0.95\linewidth] {
                    \textbf{\color{CapgeminiBlue}\small Ce que j'apporte à vos équipes :}
                    \vspace{0.12cm}
                    \begin{itemize}\setlength{\itemsep}{3.5pt}\scriptsize
                        \item \textbf{Double compétence Dév \& Test} : Capacité à lire et comprendre le code pour écrire des tests unitaires et d'API pertinents.
                        \item \textbf{Rigueur éprouvée sur le terrain} : Expérience réussie en milieu industriel chez OCP Safi (222 tests, SonarQube à 86.8\%).
                        \item \textbf{Pratique agile naturelle} : Habitué aux rituels Scrum, à la communication claire et au travail d'équipe.
                        \item \textbf{Motivation forte pour la QA} : Envie d'apprendre des méthodologies de la Testing Factory de Capgemini.
                    \end{itemize}
                };
            \end{tikzpicture}
        \end{column}

        % Colonne 2 : Modalités pratiques & Contact
        \begin{column}{0.42\textwidth}
            \begin{tikzpicture}
                \node[card, fill=CapgeminiBlue!6, draw=CapgeminiBlue!30, text width=0.92\linewidth] {
                    \centering
                    \textbf{\color{CapgeminiBlue}\normalsize Stage PFE 2027}\\[0.1cm]
                    \scriptsize
                    \textbf{Durée :} 6 mois (Plein temps)\\[0.05cm]
                    \textbf{Démarrage :} \textbf{25 Janvier 2027}\\[0.05cm]
                    \textbf{Mobilité :} Casablanca, Rabat, International\\[0.1cm]
                    {\color{CapgeminiBlue!40}\rule{2.5cm}{0.5pt}}\\[0.1cm]
                    \textbf{École :} EMSI Marrakech\\[0.05cm]
                    \textbf{Profil :} Ingénieur QA / Testeur Logiciel
                };
            \end{tikzpicture}
            \vspace{0.15cm}

            \begin{tikzpicture}
                \node[card, text width=0.92\linewidth] {
                    \centering
                    \textbf{\scriptsize Rafik AIT ICHOU}\\[0.04cm]
                    \tiny\color{TextSecondary}
                    \href{mailto:rafikaitichou@gmail.com}{rafikaitichou@gmail.com} \ \textbullet\ \ +212 600-607637\\[0.04cm]
                    \href{https://aitichou.com}{aitichou.com} \ \textbullet\ \ \href{https://linkedin.com/in/rafik-ait-ichou}{linkedin.com/in/rafik-ait-ichou}
                };
            \end{tikzpicture}
        \end{column}
    \end{columns}

    \vspace{0.25cm}
    \centering
    {\small\color{CapgeminiBlue}\bfseries Merci pour votre attention — Je suis ravi d'échanger avec vous !}
\end{frame}

\end{document}
